%%%%

%%%   Самарский государственный технический университет

%%%   Кафедра прикладной математики и информатики

%%%

%%%   Преамбула для статьи в журнал "Вестник Самарского государственного

%%%   технического университета. Серия: «Физико-математические науки»"

%%%   Ориентирована под MikTEx 2.4 for Win32

%%%

%%%   Хотя сам журнал собирается под GNU/Linux на дистрибутиве TeXLive



\documentclass[11pt,a4paper,twoside]{article}



\usepackage[cp1251]{inputenc}

\usepackage[T2A]{fontenc}

\usepackage[english,russian]{babel}

\usepackage{samgtubib}

\usepackage[dvips]{graphicx}

\usepackage[multiple]{footmisc}



\usepackage{samgtu}

\renewcommand{\VestnikYear}{2009} % Год издания Вестника

\renewcommand{\VestnikNumber}{2\,(19)} % Номер Вестника

\renewcommand{\VestnikMonth}{Ноябрь} % Месяц выхода Вестника

\renewcommand{\VestnikData}{01.11.09} % Дата подписания Вестника в печать

\renewcommand{\VestnikUPL}{26,64} % Усл. печ. л.

\renewcommand{\VestnikUIL}{25,73} % Уч.-изд. л.

\renewcommand{\VestnikRegNumber}{479/09} % Регистрационный номер

\renewcommand{\VestnikOrder}{994} % Номер заказа






%
%\begin{document}

%
\renewcommand{\baselinestretch}{.95}

\renewcommand{\firstpage}{83}
\renewcommand{\lastpage}{88}

\udk{539.3} \msc{74B05; 74K20}

\title{Упругие напряжения в плоскости с тонкостенным кусочно"=однородным включением при наличии сосредоточенных сил}
{Упругие напряжения в плоскости с тонкостенным кусочно-однородным включением \dots}% Заголовок, который идёт в верхний колонтитул

\author{В.\,В.~Сильвестров, А.\,В.~Смирнов}
{С\,и\,л\,ь\,в\,е\,с\,т\,р\,о\,в~В.\,В., С\,м\,и\,р\,н\,о\,в~А.\,В.}

\institute{Российский государственный университет нефти и~газа им.~И.\,М.~Губкина, \\ 
199991, Москва, Ленинский пр., 65. }


\email{E-mails}{v-silvestrov@yandex.ru, smirnov09al@gmail.com}

\otherinf{\fullauthor{Василий Васильевич Сильвестров}\regalia{д.ф.-м.н.}\position{профессор} \dept{каф. высшей математики}
\fullauthor{Александр Валериянович Смирнов} \position{аспирант} \dept{каф. высшей математики}}

\keywords{упругие пластины, включение, сосредоточенные силы, напряжения}

\workabstract{Рассматривается упругая пластина, составленная из двух полубесконечных пластин, соединенных друг с другом через тонкое упругое
кусочно"=однородное включение. Изучается напряжённое состояние, порожденное двумя сосредоточенными силами, приложенными к
пластинам. Находятся напряжения на линиях контакта, изучается их поведение вблизи точки изменения жёсткости включения.}

\engtitle{Elastic stresses in a~plane with a~thin piecewise-homogeneous inclusion}

\engauthor{V.\,V.~Silvestrov, A.\,V.~Smirnov,}{S\,i\,l\,v\,e\,s\,t\,r\,o\,v~V.\,V., S\,m\,i\,r\,n\,o\,v~A.\,V.}

\enginstitute{Gubkin Russian State University of Oil and Gas, \\ 65, Leninskiy pr., Moscow, 199991, Russia.}

\otherenginfm{\fullauthor{Vasily V. Silvestrov} \regalia{Dr. Sci. (Phys. \& Math.)} \position{Professor}
\dept{Dept. of Higher Mathematics}
\fullauthor{Aleksandr V. Smirnov}
\position{Postgraduate Student} \dept{Dept. of Higher Mathematics}}

\engabstract{The elastic plate composed of two semi-infinite plates connected one with other through a thin elastic piecewise-homogeneous inclusion is considered. Stress state in the plates under the action applied to them concentrated forces is studied. The stresses on the lines of contact are found and their behavior near point of changing the inclusion rigidity is studied.}

\engkeywords{elastic plates, inclusion, stresses, concentrated forces}

\date{24/I/2011}
\revisiondate{27/III/2011}



\maketitle



Пусть две полубесконечные однородные пластины с различными упругими характеристиками соединены между собой через бесконечное
тонкое упругое включение с кусочно"=постоянным модулем упругости (рис. \ref{silsmi:pic1}). Средние сечения пластин лежат в
плоскости комплексной переменной $z=x+iy$ так, что ось включения, расположенного между ними, совпадает с действительной осью.
Обозначим толщины, коэффициенты Пуассона, модули сдвига пластин через $h_1$, $\nu_1$, $E_1$ и $h_2$, $\nu_2$, $E_2$
соответственно (здесь и далее индекс 1 соответствует верхней пластине, индекс 2 "--- нижней пластине), модуль упругости включения
"--- через $E_k$ при $(-1)^{k-1}x>0$ ($k=1,\,2$), а площадь поперечного сечения включения "--- через $S$. К пластинам в точках
$z_{1,2}=\pm i d$ приложены сосредоточенные силы $P_1+i Q_1$ и $P_2+i Q_2$.

{\sl Требуется определить касательные и нормальные напряжения вдоль линий контакта пластин и включения.}

\begin{figure}[h!] 
\centering \includegraphics[scale=1]{./00PIC/fsilvestrov1}
\caption{Составная плоскость с кусочно-однородным включением}\label{silsmi:pic1}
\end{figure}


Пластины находятся в обобщённом плоском напряжённом состоянии, а включение "--- в одноосном напряжённом состоянии. При этом
из"=за различия упругих характеристик пластин происходит не только растяжение, но и изгиб включения, учет которого значительно
усложняет процесс решения. Здесь, как и в работе \cite{silsmi:bib:pmm}, предполагается, что включение лишено изгибной жёсткости,
т.\,е. совершенно не сопротивляется изгибу. В этом случае из условия равновесия части включения на интервале $(x; x+\Delta x)$,
переходя к пределу при $\Delta x\to 0$, получим уравнения
\begin{equation}
\begin{aligned}
&E_k A u''(x) + h_1 \tau_{xy}^+ (x) - h_2 \tau_{xy}^- (x)=0,\\
h_1 \sigma_y^+(x)&-h_2\sigma_y^-(x)=0, \quad (-1)^{k-1}x>0,\quad k=1,\,2, \label{silsmi:eq1}
\end{aligned}
\end{equation}
где $\sigma_y^\pm$ и $\tau_{xy}^\pm$ "--- искомые нормальные и касательные контактные напряжения в верхней и нижней пластине. В
силу малости толщины включения на его гранях должно выполняться условие равенства горизонтальных и вертикальных смещений точек
пластин \cite{silsmi:bib:chobanan}
\begin{equation}
\label{silsmi:eq2}
(u^++iv^+)(x)=(u^-+iv^-)(x),\quad x\in\mathbb{R}.
\end{equation}

Искомые напряжения и деформации в пластинах выражаются по формулам Колосова"--~Мусхелишвили \cite{silsmi:bib:mushel}:
\begin{gather*}
\sigma_y(z)-i\tau_{xy}(z)=\Phi_k(z)-\Phi_k(\overline z)+(z-\overline z)\overline{\Phi'_k(z)},\\
2\mu_k(u+i v)'(z)=\kappa_k\Phi_k(z)+\Phi_k(\overline z)-(z-\overline z)\overline{\Phi'_k(z)},\quad (-1)^{k-1}\Im z>0,\, k=1,\,2,
\end{gather*}
где $\kappa_k=(3-\nu_k)/(1+\nu_k)$, а $\Phi_k(z)$ "--- исчезающие на бесконечности мероморфные функции с линией разрыва по
действительной оси, причём функция $\Phi_k(z)$ может иметь простой полюс в точке $z_k$ и полюс второго порядка в точке $z_{3-k}$
($k=1,\,2$). Из этих формул и равенства (\ref{silsmi:eq2}) следует
\begin{equation}
\label{silsmi:Phi12link} \mu_2\kappa_1 \Phi_1^+(x)-\mu_1\Phi_2^+(x)=\mu_1\kappa_2\Phi_2^-(x)-\mu_2\Phi_1^-(x),\quad
x\in\mathbb{R}.
\end{equation}
Из свойств функций $\Phi_{1,2}(z)$ и краевого условия (\ref{silsmi:Phi12link}) вытекает аналитичность функции
\begin{equation}
\label{silsmi:R0def}
R_0(z)=\begin{cases}
\mu_2\kappa_1\Phi_1(z)-\mu_1\Phi_2(z), \quad \Im z>0,\\
\mu_1\kappa_2\Phi_2(z)-\mu_2\Phi_1(z), \quad \Im z<0
\end{cases}
\end{equation}
во всей комплексной плоскости за исключением точек $z_{1,2}$, в которых она может иметь полюсы второго порядка, причём на
бесконечности функция $R_0(z)$ обращается в ноль. Следовательно, она имеет вид:
$$
R_0(z)=\frac{A_0}{z-id}+\frac{B_0z+C_0}{(z-id)^2}+\frac{D_0}{z+id}+\frac{E_0z+F_0}{(z+id)^2}.
$$
Постоянные $A_0$, $B_0$, $C_0$, $D_0$, $E_0$, $F_0$ определяются из известных представлений комплексных потенциалов вблизи точек
приложения сосредоточенных сил \cite{silsmi:bib:mushel}:
$$
\begin{aligned}
\Phi_1(z)=&-\frac{P_1+iQ_1}{2\pi(1+\kappa_1)}\frac1{z-id}+\Phi_{01}(z),& &\Im z>0; \\
\Phi_1(z)=&\frac{(1-\kappa_1)P_1-i(1+\kappa_1)Q_1}{2\pi(1+\kappa_1)}\frac1{z+id}-  & & \\
&\qquad  -\frac{P_1-iQ_1}{2\pi(1+\kappa_1)}\frac{z-id}{(z+id)^2}+\Phi_{01}(z),& &\Im z<0;\\
\Phi_2(z)=&\frac{(1-\kappa_2)P_2-i(1+\kappa_2)Q_2}{2\pi(1+\kappa_2)}\frac1{z-id}- & &\\&\qquad  -\frac{P_2-iQ_2}{2\pi(1+\kappa_2)}\frac{z+id}{(z-id)^2}+\Phi_{02}(z),& &\Im z>0;\\
\Phi_2(z)=&-\frac{P_2+iQ_2}{2\pi(1+\kappa_2)}\frac1{z+id}+\Phi_{02}(z),& &\Im z<0;
\end{aligned}
$$
где $\Phi_{01}(z)$ и $\Phi_{02}(z)$ аналитичны в окрестностях точек $z_{1,2}$.

На основании формул Колосова"--~Мусхелишвили и равенства (\ref{silsmi:R0def}) условия (\ref{silsmi:eq1}) преобразуются в краевое
условие задачи Римана:
\begin{gather}
\label{silsmi:eqRiemann}
(1+\delta\kappa_1)\Phi_1^+(x)-(1+\delta\kappa_2^{-1})\Phi_1^-(x)=i g(x),\quad x\in\mathbb{R}\backslash\{0\},\\
g(x)=\frac{E_kA}{2\mu_1h_1}\frac{d}{dx}\Re\left[\kappa_1\Phi_1^+(x)+\Phi_1^-(x)\right]-i
\frac{h_2}{h_1}\frac{1+\kappa_2^{-1}}{\mu_1}R_0(x),\quad \delta=\frac{h_2\mu_2}{h_1\mu_1}. \label{silsmi:gdef}
\end{gather}
Функция $g(x)$ непрерывна по Гельдеру на всей действительной оси, кроме точки $x=0$, где она может обращаться в бесконечность
порядка меньше 1, и исчезает на бесконечности не медленнее $1/x$. Следовательно, решение задачи Римана имеет вид
\begin{equation}
\label{silsmi:solRiemann}
\begin{aligned}
&\Phi_1(z)=\frac{A_1}{z-id}+\frac1{2\pi(1+\delta\kappa_1)}\int _{-\infty}^{+\infty}\frac{g(t)dt}{t-z},\quad \Im z>0,\\
&\Phi_1(z)=\frac{B_1}{z+id}+\frac{C_1z+D_1}{(z+id)^2}+\frac{\kappa_2}{2\pi(\kappa_2+\delta)}\int _{-\infty}^{+\infty}\frac{g(t)dt}{t-z},\quad \Im z<0.
\end{aligned}
\end{equation}
Постоянные $A_1$, $B_1$, $C_1$, $D_1$ снова определяются из приведенных представлений комплексных потенциалов.

Из равенства (\ref{silsmi:gdef}) находим $\Im g(x)=-\delta\mu_2^{-1}(1+\kappa_2^{-1})\Re R_0(x)$. Для нахождения $\Re g(x)$
подставим предельные значения интегралов (\ref{silsmi:solRiemann}) в равенство (\ref{silsmi:gdef}) и выделим действительные части
справа и слева. Вводя новую функцию $f'(x)=\Re g(x)$, получим систему интегро"=дифференциальных уравнений Прандтля
\begin{equation}
\label{silsmi:eqPrandtl}
\begin{aligned}
&f(x)-\frac{a_1}{\pi}\int _{-\infty}^{+\infty}\frac{f'(x)dt}{t-x}=E_1 r_1(x)-\int _x^{+\infty}r_2(t)dt,& & x>0,\\
&f(x)-\frac{a_2}{\pi}\int _{-\infty}^{+\infty}\frac{f'(x)dt}{t-x}=E_2 r_1(x)+\int _{-\infty}^{x}r_2(t)dt,& & x<0,
\end{aligned}
\end{equation}
где
\begin{gather*}
r_1(x)=\frac{A}{2\mu_1h_1}\Re\left[R_1(x)+\frac12\left(\frac{\kappa_1}{1+\delta\kappa_1}-\frac{\kappa_2}{\kappa_2+\delta}\right)\frac{h_2}{h_1}\frac{1+\kappa_2^{-1}}{\mu_1}R_0(x)\right],\\
r_2(x)=\frac{h_2}{h_1}\frac{1+\kappa_2^{-1}}{\mu_1}\Im R_0(x),\quad R_1(x)=\frac{\kappa_1A_1}{x-id}+\frac{B_1}{x+id}+\frac{C_1x+D_1}{(x+id)^2},\\
a_k=\frac{E_kA}{4\mu_1h_1}\left(\frac{\kappa_1}{1+\delta\kappa_1}-\frac{\kappa_2}{\kappa_2+\delta}\right).
\end{gather*}

Система уравнений (\ref{silsmi:eqPrandtl}) допускает решение методом, подробно описанным в работе \cite{silsmi:bib:pmm}.
Посредством интегрального преобразования Меллина она сводится к системе двух разностных уравнений, которая с помощью выбранной
специальным образом диагонализации и последующего конформного отображения единичной полосы на расширенную комплексную плоскость
преобразуется в краевую задачу сопряжения на римановой поверхности. Ввиду громоздкости решение здесь не приводится. Найденные на
основе этого решения контактные напряжения удовлетворяют условию Гельдера во всех точках действительной оси, кроме точки $x=0$.
При этом касательные напряжения имеют логарифмическую особенность в окрестности нуля, а нормальные напряжения принимают различные
конечные значения при стремлении к нулю справа и слева.

\begin{figure}[h!]
\centering \includegraphics[width=\textwidth]{./00PIC/fsilvestrov2}

\footnotesize\sl a \hspace{.48\textwidth }б

\caption{Контактные напряжения вблизи нуля: \footnotesize  $P_1=P_2=P_0$ ({\sl а\/}); $P_1=-P_2=P_0$ ({\sl б\/})}\label{silsmi:pic2}
\end{figure}



На рис.~\ref{silsmi:pic2} представлены графики контактных напряжений при горизонтальной нагрузке ($Q_1=Q_2=0$) в двух случаях
$P_1=P_2$ и $P_1=-P_2$. На рис.~\ref{silsmi:pic3} изображены контактные напряжения при различных значениях величины $d$ при
внешней нагрузке $P_1=-P_2=P_0$. Из рисунков видно, что при уменьшении расстояния между точками приложения сосредоточенных сил и
включением контактные напряжения испытывают сильное возмущение вблизи нуля, причём на некоторых участках нормальные напряжения по
величине превышают касательные, тогда как при больших $d$ нормальные напряжения значительно меньше касательных.

\thanks{Работа выполнена при поддержке РФФИ (проект \rm \No~10--01-00103--a).}

\newpage

\begin{figure}[h!]
\centering \includegraphics[width=\textwidth]{./00PIC/fsilvestrov3}

\caption{Касательные и нормальные контактные напряжения при различных значениях расстояния $d$: \footnotesize 1 "--- $d=h_0$, 
2 "--- $d=2h_0$, 3 "--- $d=4h_0$, 4 "--- $d=8h_0$}\label{silsmi:pic3}
\end{figure}

\begin{thebibliography}{9}

\RBibitem{silsmi:bib:pmm}
\paper Интегро-дифференциальное уравнение Прандтля и контактная задача для кусочно--однородной пластины
\by Сильвестров~В.\,В., Смирнов~А.\,В.
\jour ПММ
\yr 2010
\vol 74
\issue 6
\pages 951--968
\transl англ. пер.:~
\paper The Prandtl integrodifferential equation and the contact problem for a piecewise homogeneous plate
\by Sil'vestrov~V.\,V., Smirnov~A.\,V.
\jour J.~Appl. Math. Mech. 
\vol 74
\issue 6
\pages 679--691 
\yr 2010

\RBibitem{silsmi:bib:chobanan}
\paper Плоское деформированное состояние упругого тела с тонкостенным гибким включением
\by Чобанян~К.\,С.,Хачикян~А.\,С.
\jour Изв. АН АрмССР. Сер. Механика
\yr 1967
\issue 6
\pages 19--29
\misctransl
\by Chobanyan~K.\,S., Khachikyan~A.\,S.
\paper Plane deformation state of an elastic body with a thin walled flexible inclusion
\yr 1967
\issue 6
\pages 19--29
\jour Izv. AN ArmSSR. Ser. Mehanika

\RBibitem{silsmi:bib:mushel}
\by Мусхелишвили~Н.\,И.
\book Некоторые основные задачи математической теории упругости
\publaddr М.
\publ Наука
\yr 1966
\totalpages 707
\transl
англ. пер.:~
\by Muskhelishvili~N.\,I.
\book Some basic problems of the mathematical theory of elasticity. Basic equations, the plane theory of elasticity, torsion and bending
\yr 1977
\publ Noordhoff International Publishing
\publaddr Leyden
\totalpages 768

\end{thebibliography}

\noindent
\hskip6.5cm \thedate \par\noindent
\hskip6.5cm\therevdate\par

\newpage

\chead[\fancyplain{}{\scriptsize \sl S\,i\,l\,v\,e\,s\,t\,r\,o\,v~V.\,V., S\,m\,i\,r\,n\,o\,v~A.\,V.}]{\fancyplain{}{\scriptsize \sl  Elastic stresses in a~plane with a~thin piecewise-homogeneous inclusion}}


\makeengtitlem

\renewcommand{\baselinestretch}{0.9}

\newpage

%\end{document}
